%%%%%%%%%%%%%%%%%%%%%%%%%%%%%%%%%%%%%%%%%%%%%%%%%%%%%%%%%%%%%%%%%%%%%%%%%%%%%%%%%%
%  RAG PIPELINE: END-TO-END IMPLEMENTATION
%
%  Tools : LangChain v0.1 (v1 API), DocstringLoader, HuggingFace Embeddings,
%           ChromaDB, HuggingFaceHub LLM, RAGAS
%  Data  : local faq.txt
%  Goal  : academic walk-through of every pipeline step with runnable code
%%%%%%%%%%%%%%%%%%%%%%%%%%%%%%%%%%%%%%%%%%%%%%%%%%%%%%%%%%%%%%%%%%%%%%%%%%%%%%%%%%

\begin{frame}[fragile]\frametitle{}
\begin{center}
{\Large RAG Pipeline: End-to-End Implementation}
\end{center}
\end{frame}


%%%%%%%%%%%%%%%%%%%%%%%%%%%%%%%%%%%%%%%%%%%%%%%%%%%%%%%%%%%%%%%%%%%%%%%%%%%%%%%%%%
\begin{frame}[fragile]\frametitle{Implementation Overview}

\textbf{Goal:} Build a question-answering system over a local \texttt{faq.txt} file and measure its quality with RAGAS: using only open-source components.

\vspace{0.4em}
\begin{center}
\small
\texttt{faq.txt}
$\xrightarrow{\text{1. Load}}$
\texttt{Documents}
$\xrightarrow{\text{2. Chunk}}$
\texttt{Chunks}
$\xrightarrow{\text{3. Embed}}$
\texttt{Vectors}
$\xrightarrow{\text{4. Store}}$
\texttt{ChromaDB}
$\xrightarrow{\text{5. Retrieve}}$
\texttt{Context}
$\xrightarrow{\text{6. Generate}}$
\texttt{Answer}
$\xrightarrow{\text{7. Evaluate}}$
\texttt{RAGAS Scores}
\end{center}

% \vspace{0.4em}
% \begin{center}
% \small
% \begin{tabular}{lll}
% \hline
% \textbf{Step} & \textbf{Tool} & \textbf{LangChain class} \\
% \hline
% Load \& Parse    & DocstringLoader         & \texttt{langchain.document\_loaders} \\
% Chunk            & Recursive char splitter & \texttt{langchain.text\_splitter} \\
% Embed            & HuggingFace MiniLM      & \texttt{langchain.embeddings} \\
% Store            & ChromaDB                & \texttt{langchain.vectorstores} \\
% Retrieve         & Similarity search       & \texttt{VectorStoreRetriever} \\
% Generate         & HuggingFaceHub LLM      & \texttt{langchain.llms} \\
% Chain            & RetrievalQA             & \texttt{langchain.chains} \\
% Evaluate         & RAGAS                   & \texttt{ragas} \\
% \hline
% \end{tabular}
% \end{center}

\end{frame}


%%%%%%%%%%%%%%%%%%%%%%%%%%%%%%%%%%%%%%%%%%%%%%%%%%%%%%%%%%%%%%%%%%%%%%%%%%%%%%%%%%
\begin{frame}[fragile]\frametitle{The Data: faq.txt}

\begin{columns}
  \begin{column}[T]{0.6\linewidth}
    \begin{itemize}
      \item A plain-text FAQ file: the simplest possible real-world knowledge source.
      \item Each entry follows a \textbf{Q:} / \textbf{A:} format; entries are separated by a blank line.
      \item Plain text is ideal for an academic demo: no PDF parsing complexity, no OCR errors, fully reproducible.
      \item The same pattern scales to any \texttt{.txt}, \texttt{.md}, or \texttt{.rst} knowledge base.
      \item DocstringLoader will read the entire file into a single LangChain \texttt{Document} object; the splitter then breaks it into per-entry chunks.
    \end{itemize}
  \end{column}
  \begin{column}[T]{0.4\linewidth}
\begin{lstlisting}[language=bash, basicstyle=\tiny]
# faq.txt  (excerpt)

Q: What is RAG?
A: Retrieval-Augmented Generation (RAG)
   combines a retrieval system with a
   language model. The retriever fetches
   relevant documents; the LLM generates
   an answer grounded in those documents.

Q: Why use RAG instead of fine-tuning?
A: RAG keeps knowledge outside the model
   weights, so it can be updated without
   retraining. Fine-tuning is expensive
   and risks forgetting prior knowledge.

Q: What is a vector database?
A: A database optimised for storing and
   searching high-dimensional embedding
   vectors using approximate nearest-
   neighbour (ANN) algorithms.

Q: What is chunking?
A: Splitting documents into smaller pieces
   so each piece fits in the LLM context
   window and carries a focused meaning.
\end{lstlisting}
  \end{column}
\end{columns}

\end{frame}


%%%%%%%%%%%%%%%%%%%%%%%%%%%%%%%%%%%%%%%%%%%%%%%%%%%%%%%%%%%%%%%%%%%%%%%%%%%%%%%%%%
\begin{frame}[fragile]\frametitle{Dependencies and Setup}

\begin{columns}
  \begin{column}[T]{0.6\linewidth}
    \begin{itemize}
      \item \textbf{langchain} v0.1: the v1 stable API used throughout this demo.
      \item \textbf{sentence-transformers}: local embedding model; no API key needed.
      \item \textbf{chromadb}: lightweight local vector store; persists to disk.
      \item \textbf{huggingface\_hub}: access to open-source LLMs via Inference API.
      \item \textbf{ragas}: reference-free RAG evaluation framework.
      \item \textbf{datasets}: required by RAGAS to build the evaluation dataset object.
      \item Set \texttt{HUGGINGFACEHUB\_API\_TOKEN} in your environment (free tier is sufficient for demos).
    \end{itemize}
  \end{column}
  \begin{column}[T]{0.4\linewidth}
\begin{lstlisting}[language=bash, basicstyle=\tiny]
# Create a virtual environment (recommended)
conda create -n rag-demo
conda activate rag-demo

# Install dependencies
pip install langchain==0.1.20
pip install langchain-community==0.0.38
pip install sentence-transformers==2.7.0
pip install chromadb==0.5.0
pip install huggingface_hub==0.23.0
pip install ragas==0.1.9
pip install datasets==2.19.0

# Set your HuggingFace token
export HUGGINGFACEHUB_API_TOKEN="hf_..."
\end{lstlisting}
  \end{column}
\end{columns}

\end{frame}


%%%%%%%%%%%%%%%%%%%%%%%%%%%%%%%%%%%%%%%%%%%%%%%%%%%%%%%%%%%%%%%%%%%%%%%%%%%%%%%%%%
\begin{frame}[fragile]\frametitle{Step 1: Load and Parse with DocstringLoader}

\begin{columns}
  \begin{column}[T]{0.6\linewidth}
    \begin{itemize}
      \item \textbf{DocstringLoader} reads a text file and returns a list of LangChain \texttt{Document} objects, each carrying \texttt{page\_content} and \texttt{metadata}.
      \item Unlike the bare \texttt{TextLoader}, DocstringLoader preserves \textbf{paragraph structure}: blank-line boundaries become natural document boundaries.
      \item Each \texttt{Document} automatically receives \texttt{source} metadata pointing to the file path, useful for citations in the final answer.
      \item Printing the first document confirms the content was loaded correctly before proceeding.
    \end{itemize}
  \end{column}
  \begin{column}[T]{0.4\linewidth}
\begin{lstlisting}[language=python, basicstyle=\tiny]
from langchain.document_loaders import TextLoader


loader = TextLoader(
    file_path="faq.txt",
    encoding="utf-8"
)
documents = loader.load()

print(f"Loaded {len(documents)} document(s)")
print("--- First document:")
print(documents[0].page_content[:300])
print("Metadata:", documents[0].metadata)

# Expected output:
# Loaded 1 document(s)
#: First document:
# Q: What is RAG?
# A: Retrieval-Augmented Generation ...
# Metadata: {'source': 'faq.txt'}
\end{lstlisting}
  \end{column}
\end{columns}

{\tiny (Ref: https://python.langchain.com/v0.1/docs/modules/data\_connection/document\_loaders/)}

\end{frame}


%%%%%%%%%%%%%%%%%%%%%%%%%%%%%%%%%%%%%%%%%%%%%%%%%%%%%%%%%%%%%%%%%%%%%%%%%%%%%%%%%%
\begin{frame}[fragile]\frametitle{Step 2: Chunking the Document}

\begin{columns}
  \begin{column}[T]{0.6\linewidth}
    \begin{itemize}
      \item \textbf{RecursiveCharacterTextSplitter} is the recommended default for plain text: it tries paragraph, then sentence, then word boundaries before falling back to characters.
      \item \texttt{chunk\_size=300}: each chunk covers roughly one FAQ entry (question + answer).
      \item \texttt{chunk\_overlap=50}: a 50-character tail is repeated in the next chunk, preventing an answer from being split across a boundary.
      \item \texttt{separators} list prioritises blank lines (\texttt{\textbackslash n\textbackslash n}) so each Q\&A pair stays together.
      \item Inspect chunk count and the first chunk to verify boundaries are sensible.
    \end{itemize}
  \end{column}
  \begin{column}[T]{0.4\linewidth}
\begin{lstlisting}[language=python, basicstyle=\tiny]
from langchain.text_splitter import  RecursiveCharacterTextSplitter

splitter = RecursiveCharacterTextSplitter(
    chunk_size=300,
    chunk_overlap=50,
    separators=["\n\n", "\n", " ", ""]
)

chunks = splitter.split_documents(documents)

print(f"Total chunks: {len(chunks)}")
print("\n--- Chunk 0:")
print(chunks[0].page_content)
print("\n--- Chunk 1:")
print(chunks[1].page_content)

# Each chunk should contain one Q&A pair.
# chunk_overlap ensures no answer is lost
# at a boundary.
\end{lstlisting}
  \end{column}
\end{columns}

{\tiny (Ref: https://python.langchain.com/v0.1/docs/modules/data\_connection/text\_splitting/)}

\end{frame}


%%%%%%%%%%%%%%%%%%%%%%%%%%%%%%%%%%%%%%%%%%%%%%%%%%%%%%%%%%%%%%%%%%%%%%%%%%%%%%%%%%
\begin{frame}[fragile]\frametitle{Step 3: Embed and Store in ChromaDB}

\begin{columns}
  \begin{column}[T]{0.6\linewidth}
    \begin{itemize}
      \item \textbf{HuggingFaceEmbeddings} wraps the \texttt{sentence-transformers} library. \texttt{all-MiniLM-L6-v2} maps text to 384-dimensional vectors: lightweight and fast on CPU.
      \item \textbf{Chroma.from\_documents()} embeds every chunk and stores vectors plus original text on disk under \texttt{./chroma\_db}.
      \item \texttt{persist\_directory} enables loading the index in future sessions without re-embedding.
      \item Printing the collection count confirms all chunks were indexed successfully.
    \end{itemize}
  \end{column}
  \begin{column}[T]{0.4\linewidth}
\begin{lstlisting}[language=python, basicstyle=\tiny]
from langchain.embeddings import HuggingFaceEmbeddings
from langchain.vectorstores import Chroma

# Local embedding model - no API key needed
embedding_model = HuggingFaceEmbeddings(
    model_name="all-MiniLM-L6-v2",
    model_kwargs={"device": "cpu"}
)

# Build and persist the vector store
vectorstore = Chroma.from_documents(
    documents=chunks,
    embedding=embedding_model,
    persist_directory="./chroma_db"
)
vectorstore.persist()

count = vectorstore._collection.count()
print(f"Indexed {count} chunks in ChromaDB")
\end{lstlisting}
  \end{column}
\end{columns}

{\tiny (Ref: https://python.langchain.com/v0.1/docs/integrations/vectorstores/chroma/)}

\end{frame}


%%%%%%%%%%%%%%%%%%%%%%%%%%%%%%%%%%%%%%%%%%%%%%%%%%%%%%%%%%%%%%%%%%%%%%%%%%%%%%%%%%
\begin{frame}[fragile]\frametitle{Step 4: Retriever and Prompt Template}

\begin{columns}
  \begin{column}[T]{0.6\linewidth}
    \begin{itemize}
      \item \texttt{as\_retriever(search\_kwargs=\{"k": 3\})} returns a \texttt{VectorStoreRetriever} that runs cosine similarity search and returns the top-3 chunks.
      \item The \textbf{PromptTemplate} injects both \texttt{\{context\}} (retrieved chunks) and \texttt{\{question\}} (user query) into a single formatted string.
      \item The grounding instruction: ``\textit{use only the provided context}'': is the key safeguard against hallucination. If the answer is absent from context, the model is instructed to say so.
      \item \texttt{input\_variables} must match exactly the placeholder names used in the template string.
    \end{itemize}
  \end{column}
  \begin{column}[T]{0.4\linewidth}
\begin{lstlisting}[language=python, basicstyle=\tiny]
from langchain.prompts import PromptTemplate

# Retriever: return top-3 similar chunks
retriever = vectorstore.as_retriever(
    search_type="similarity",
    search_kwargs={"k": 3}
)

# Grounded prompt template
template = """You are a helpful assistant.
Answer the question using ONLY the context
provided below. If the answer is not in
the context, say "I don't know."

Context:
{context}

Question: {question}

Answer:"""

prompt = PromptTemplate(
    input_variables=["context", "question"],
    template=template
)
\end{lstlisting}
  \end{column}
\end{columns}

\end{frame}


%%%%%%%%%%%%%%%%%%%%%%%%%%%%%%%%%%%%%%%%%%%%%%%%%%%%%%%%%%%%%%%%%%%%%%%%%%%%%%%%%%
\begin{frame}[fragile]\frametitle{Step 5: LLM and RetrievalQA Chain}

\begin{columns}
  \begin{column}[T]{0.6\linewidth}
    \begin{itemize}
      \item \textbf{HuggingFaceHub} calls the HuggingFace Inference API: no local GPU needed. \texttt{google/flan-t5-large} is a compact instruction-following model well-suited to Q\&A.
      \item \texttt{max\_new\_tokens=256} caps response length; \texttt{temperature=0} makes output deterministic: important for reproducible academic experiments.
      \item \textbf{RetrievalQA.from\_chain\_type()} wires the retriever, prompt, and LLM into a single callable object.
      \item \texttt{return\_source\_documents=True} exposes which chunks were used, enabling manual inspection and RAGAS evaluation.
    \end{itemize}
  \end{column}
  \begin{column}[T]{0.4\linewidth}
\begin{lstlisting}[language=python, basicstyle=\tiny]
from langchain.llms import HuggingFaceHub
from langchain.chains import RetrievalQA

# Open-source LLM via HuggingFace Inference API
llm = HuggingFaceHub(
    repo_id="google/flan-t5-large",
    model_kwargs={
        "temperature": 0,
        "max_new_tokens": 256
    }
)

# Wire retriever + prompt + LLM into one chain
qa_chain = RetrievalQA.from_chain_type(
    llm=llm,
    chain_type="stuff",
    retriever=retriever,
    return_source_documents=True,
    chain_type_kwargs={"prompt": prompt}
)
\end{lstlisting}
  \end{column}
\end{columns}

{\tiny (Ref: https://python.langchain.com/v0.1/docs/modules/chains/)}

\end{frame}


%%%%%%%%%%%%%%%%%%%%%%%%%%%%%%%%%%%%%%%%%%%%%%%%%%%%%%%%%%%%%%%%%%%%%%%%%%%%%%%%%%
\begin{frame}[fragile]\frametitle{Step 6: Query the Pipeline}

\begin{columns}
  \begin{column}[T]{0.6\linewidth}
    \begin{itemize}
      \item Calling \texttt{qa\_chain()} with a query string triggers the full pipeline in one step: embed query $\rightarrow$ retrieve top-3 chunks $\rightarrow$ format prompt $\rightarrow$ LLM generates answer.
      \item The result dict contains \texttt{"result"} (the answer string) and \texttt{"source\_documents"} (the retrieved \texttt{Document} objects).
      \item Printing source documents lets you verify the retriever fetched relevant chunks: the first diagnostic step when answers look wrong.
      \item Inspect \texttt{page\_content} of each source doc to confirm it contains evidence for the answer.
    \end{itemize}
  \end{column}
  \begin{column}[T]{0.4\linewidth}
\begin{lstlisting}[language=python, basicstyle=\tiny]
query = "What is RAG and why is it useful?"

result = qa_chain({"query": query})

print("Question:", query)
print("\nAnswer:", result["result"])
print("\n--- Retrieved chunks:")
for i, doc in enumerate(
        result["source_documents"], 1):
    print(f"\n[Chunk {i}]")
    print(doc.page_content[:200])

# Sample output:
# Answer: RAG combines a retrieval system
# with a language model. It is useful
# because knowledge can be updated without
# retraining the model weights.
#
# [Chunk 1]
# Q: What is RAG?
# A: Retrieval-Augmented Generation ...
\end{lstlisting}
  \end{column}
\end{columns}

\end{frame}


%%%%%%%%%%%%%%%%%%%%%%%%%%%%%%%%%%%%%%%%%%%%%%%%%%%%%%%%%%%%%%%%%%%%%%%%%%%%%%%%%%
\begin{frame}[fragile]\frametitle{Step 7: RAGAS Evaluation Dataset}

\begin{columns}
  \begin{column}[T]{0.6\linewidth}
    \begin{itemize}
      \item RAGAS requires four parallel lists for evaluation:
        \begin{itemize}
          \item \textbf{questions}: the test queries.
          \item \textbf{answers}: LLM responses from our pipeline.
          \item \textbf{contexts}: the retrieved chunk texts for each query (list of lists).
          \item \textbf{ground\_truths}: known correct answers (list of lists); needed for Recall and Correctness metrics.
        \end{itemize}
      \item Run the pipeline for each test question in a loop; collect the answer and source document texts.
      \item Wrap everything in a HuggingFace \texttt{Dataset} object: the format RAGAS \texttt{evaluate()} expects.
    \end{itemize}
  \end{column}
  \begin{column}[T]{0.4\linewidth}
\begin{lstlisting}[language=python, basicstyle=\tiny]
from datasets import Dataset

test_questions = [
    "What is RAG?",
    "Why use RAG instead of fine-tuning?",
    "What is a vector database?",
    "What is chunking in RAG?"
]
ground_truths = [
    ["RAG combines retrieval with an LLM."],
    ["RAG updates knowledge without retraining."],
    ["A DB for storing embedding vectors."],
    ["Splitting documents into smaller pieces."]
]

answers, contexts = [], []
for q in test_questions:
    res = qa_chain({"query": q})
    answers.append(res["result"])
    contexts.append([
        d.page_content
        for d in res["source_documents"]
    ])

eval_dataset = Dataset.from_dict({
    "question":     test_questions,
    "answer":       answers,
    "contexts":     contexts,
    "ground_truths": ground_truths
})
\end{lstlisting}
  \end{column}
\end{columns}

{\tiny (Ref: https://docs.ragas.io/en/stable/getstarted/evaluation/)}

\end{frame}


%%%%%%%%%%%%%%%%%%%%%%%%%%%%%%%%%%%%%%%%%%%%%%%%%%%%%%%%%%%%%%%%%%%%%%%%%%%%%%%%%%
\begin{frame}[fragile]\frametitle{Step 8: RAGAS Metrics and Results}

\begin{columns}
  \begin{column}[T]{0.65\linewidth}
    \begin{itemize}
      \item \textbf{faithfulness}: are all claims in the answer verifiable from retrieved context? Scores below 0.7 indicate hallucination.
      \item \textbf{answer\_relevancy}: does the answer address the question? Low score means the response is off-topic.
      \item \textbf{context\_recall}: did the retriever fetch all evidence present in the ground truth?
      \item \textbf{context\_precision}: what fraction of retrieved chunks was actually useful?
      \item All scores are in $[0,1]$; higher is better. For a clean FAQ, scores above 0.85 are achievable with default settings.
    \end{itemize}

    % \vspace{0.2em}
    % \begin{center}
    % \small
    % \begin{tabular}{ll}
    % \hline
    % \textbf{Metric} & \textbf{Sample} \\
    % \hline
    % Faithfulness      & 0.92 \\
    % Answer Relevancy  & 0.88 \\
    % Context Recall    & 0.87 \\
    % Context Precision & 0.90 \\
    % \hline
    % \end{tabular}
    % \end{center}
  \end{column}
  \begin{column}[T]{0.35\linewidth}
\begin{lstlisting}[language=python, basicstyle=\tiny]
from ragas import evaluate
from ragas.metrics import faithfulness, answer_relevancy, context_recall, context_precision

# RAGAS uses an LLM internally as judge.
# Default judge model: gpt-3.5-turbo.
# For a fully local setup use a local LLM:
# from ragas.llms import LangchainLLMWrapper
# ragas_llm = LangchainLLMWrapper(llm)

result = evaluate(
    dataset=eval_dataset,
    metrics=[
        faithfulness,
        answer_relevancy,
        context_recall,
        context_precision
    ]
)

# Print scores as a pandas DataFrame
df = result.to_pandas()
print(df[["question",
          "faithfulness",
          "answer_relevancy",
          "context_recall",
          "context_precision"]])
\end{lstlisting}
  \end{column}
\end{columns}

{\tiny (Ref: https://docs.ragas.io/en/stable/concepts/metrics/)}

\end{frame}

%%%%%%%%%%%%%%%%%%%%%%%%%%%%%%%%%%%%%%%%%%%%%%%%%%%%%%%%%%%%%%%%%%%%%%%%%%%%%%%%%%
\begin{frame}[fragile]\frametitle{Key design decisions and their rationale}

\begin{itemize}
  \item \textbf{\texttt{chunk\_size=300}, \texttt{k=3}:} Each FAQ entry fits in one chunk; 3 chunks give enough context without overflowing the prompt.
  \item \textbf{Grounding instruction in prompt:} The single highest-impact step to reduce hallucination: forces the LLM to stay within retrieved evidence.
  \item \textbf{\texttt{temperature=0}:} Makes every run deterministic; essential for reproducible academic experiments.
  \item \textbf{RAGAS without ground truth:} \textit{faithfulness} and \textit{answer\_relevancy} need no labels: run them in production for continuous monitoring.
  \item \textbf{Scaling up:} Replace \texttt{faq.txt} with any document set; swap \texttt{Chroma} for Qdrant or Pinecone; replace \texttt{flan-t5-large} with any LangChain-compatible LLM.
\end{itemize}

\end{frame}


% %%%%%%%%%%%%%%%%%%%%%%%%%%%%%%%%%%%%%%%%%%%%%%%%%%%%%%%%%%%%%%%%%%%%%%%%%%%%%%%%%%
% \begin{frame}[fragile]\frametitle{Complete Pipeline Summary}

% \begin{center}
% \small
% \begin{tabular}{clll}
% \hline
% \textbf{Step} & \textbf{What happens} & \textbf{Class used} & \textbf{Key parameter} \\
% \hline
% 1 & Load \texttt{faq.txt}     & \texttt{TextLoader}                       & \texttt{encoding="utf-8"} \\
% 2 & Chunk text                & \texttt{RecursiveCharacterTextSplitter}    & \texttt{chunk\_size=300} \\
% 3 & Embed chunks              & \texttt{HuggingFaceEmbeddings}            & \texttt{all-MiniLM-L6-v2} \\
% 4 & Store vectors             & \texttt{Chroma.from\_documents()}         & \texttt{persist\_directory} \\
% 5 & Retrieve top-3            & \texttt{VectorStoreRetriever}             & \texttt{k=3} \\
% 6 & Format prompt             & \texttt{PromptTemplate}                   & grounding instruction \\
% 7 & Generate answer           & \texttt{HuggingFaceHub}                   & \texttt{flan-t5-large} \\
% 8 & Evaluate                  & \texttt{ragas.evaluate()}                 & 4 metrics \\
% \hline
% \end{tabular}
% \end{center}

% \end{frame}
